\documentclass[11pt,a4paper]{article}

\usepackage[margin=2.5cm]{geometry}
\usepackage{amsmath,amssymb,bm}
\usepackage{xcolor}
\usepackage{enumitem}
\usepackage[colorlinks,linkcolor=blue,urlcolor=blue]{hyperref}

\newcommand{\vek}[1]{\bm{#1}}
\newcommand{\grad}{\vek\nabla}
\newenvironment{referee}{\par\medskip\noindent\color{black!70}\itshape}{\par\medskip}
\newcommand{\reply}{\noindent\textbf{Reply:}\ }
\newcommand{\changes}{\par\smallskip\noindent\textbf{Changes in the manuscript:}\ }

\begin{document}

\begin{center}
{\Large\bfseries Reply to the report of Referee 1}\\[0.5em]
Manuscript: \emph{Chiral soliton lattice in inhomogeneous magnetic fields}\\
T.~Brauner and R.~Radhakrishnan
\end{center}

\bigskip

We thank the referee for a careful reading of our manuscript, for the very positive assessment, and for the questions, which helped us improve the presentation and add physical content to the paper. Below we reply to all comments point by point. The referee's comments are reproduced in italics. Equation and figure numbers refer to the revised manuscript unless stated otherwise. All changes are highlighted in the accompanying marked-up version of the manuscript, with the exception of the notational change described under point~(1), which affects too many places to be usefully highlighted.

%--------------------------------------------------------------------
\section*{Point (1): notation $\vek H$}

\begin{referee}
Isn't the notation $H$ unfortunate because it usually refers to the auxiliary magnetic field in the presence of a medium? Although there is a medium here I think $H$ does not have this usual meaning, which may cause a slight, but unnecessary, confusion.
\end{referee}

\reply We agree. Our $\vek H$ was simply the external magnetic field rescaled by a constant, $\mu\vek B/(4\pi^2f_\pi^2)$, and has nothing to do with the auxiliary field of magnetostatics. We have therefore renamed it to $\vek b$ throughout the paper, including all derived quantities ($\bar{\vek b}$, $\bar b_\text{CSL}$, $\bar b_\text{cr}$, $b_0$, $b_\text{DW}$, $b_\xi$, etc.).

\changes The symbol $\vek H\to\vek b$ has been replaced everywhere in the text, equations and figure captions. Below Eq.~(3), we have added an explicit remark that $\vek b$ is merely the rescaled external magnetic field and must not be confused with the auxiliary field $\vek H$. The axis label of Fig.~1 has been updated accordingly. While doing so, we have also noticed that the stream function of the Gaussian magnetic field in Eq.~(32) was denoted by the same letter $\psi$ as the scalar potential $\psi[\vek b]$ of the gradient--solenoid decomposition~(18), although the two are unrelated. We have renamed the former to $\chi$.

%--------------------------------------------------------------------
\section*{Point (2): finite-volume phase transition and Fig.~1}

\begin{referee}
I am confused by the discussion at the end of Sec~2 and by Fig~1. Eqs~(14) and~(15) give two different classes of solutions with $\phi$ nonzero. Then there is the third possibility of $\phi=0$. Why is the transition between the former two identified with the analog of the CSL transition? Does that mean in Fig~1 there is a $\phi\neq0$ state everywhere, below and above the curve? I suspect I am misunderstanding the explanation, but since other readers might have the same problem, a clarification would be useful.
\end{referee}

\reply The referee's reading is correct: in a finite volume, the ground state is nonuniform everywhere in Fig.~1, both below and above the curve. We apologize for the insufficiently clear explanation. The key point is that in a finite volume, the natural boundary condition $\vek n\cdot\grad\phi_0=\vek n\cdot\vek b$ [Eq.~(17), formerly Eq.~(16)] is violated by $\phi=0$ for any nonzero magnetic field. Thus, $\phi=0$ is not even a stationary point of the energy functional, and the ``third possibility'' does not exist. The distinction between the vacuum and the CSL phase in a finite volume can therefore not be based on whether the pion field vanishes or not. Instead, it is based on whether the ground state contains a soliton:
\begin{itemize}[nosep]
\item The solutions of the type~(16) [formerly~(15)] satisfy $\phi_0(0)=0$ modulo $2\pi$. In weak fields, the profile stays close to the vacuum value throughout the interior and only bends in a boundary layer of thickness $\sim1/m_\pi$; for $\bar b\ll1$, linearization gives $\phi_0(\bar z)\approx\bar b\sinh\bar z/\cosh(\bar L/2)$. The energy of this deformation is a surface effect that does not contribute to the energy density as $\bar L\to\infty$. This is the finite-volume counterpart of the trivial vacuum.
\item The solutions of the type~(15) [formerly~(14)] satisfy $\phi_0(0)=\pi$ modulo $2\pi$: the center of the domain hosts the core of a soliton. This is the finite-volume analog of the single domain wall that marks the onset of the CSL phase in infinite volume.
\end{itemize}
The curve in Fig.~1 thus marks the field at which it first becomes favorable for a soliton to enter the system, not the point where $\phi_0$ becomes nonzero. In the limit $\bar L\to\infty$, the boundary layers become irrelevant and the curve approaches the infinite-volume critical field $\bar b_\text{CSL}=4/\pi$. The rapid growth of the critical field for small $\bar L$ reflects the cost of squeezing a soliton, whose natural width is of order $1/m_\pi$, into a domain of size $\lesssim1/m_\pi$. At yet stronger fields, additional solitons enter. By the reflection symmetry of the ground state, states with an odd (even) number of solitons belong to the class~(15) [(16)]. We have checked all of these statements against an independent numerical minimization of the one-dimensional energy functional.

\changes At the end of Sec.~2.2, we have added a paragraph explaining that $\phi=0$ is not a stationary point in a finite volume, and that $\phi_0\neq0$ on both sides of the curve in Fig.~1. The subsequent paragraph has been rewritten to explain the physical meaning of the two classes of solutions (the value of $\phi_0(0)$, the boundary layer, the soliton at the center) and the meaning of the curve in Fig.~1. We have also added a remark on the behavior at small $\bar L$ and on the transitions at larger fields. The caption of Fig.~1 has been expanded accordingly. We have also corrected two minor typos in this subsection: $2L\to2\bar L$ in Eq.~(14), and $\partial_z\phi_0=H\to\partial_{\bar z}\phi_0=\bar b$ in the text below Eq.~(17).

%--------------------------------------------------------------------
\section*{Point (3): physical lessons from Figs.~3 and~4}

\begin{referee}
Are there any physical lessons to be learned from the obtained configurations, say in Figs~3 and~4? I understand that the focus of the paper is mostly on developing the mathematical tools, but the reader might wonder if the behavior of the baryon number is expected or surprising etc. Some physical explanation/intuition would then also be useful to predict the qualitative behavior of configurations with more complicated inhomogeneities in the magnetic field.
\end{referee}

\reply We thank the referee for this suggestion, which prompted us to analyze the obtained configurations further. Both figures can in fact be understood analytically, and they illustrate a few general mechanisms. We have added the following discussion to the manuscript.

\begin{enumerate}[label=(\alph*)]
\item \emph{Frustration.} The anomalous term favors alignment of $\grad\phi$ with $\vek b$, but a gradient has vanishing circulation. The pion field therefore cannot follow the magnetic field wherever the field lines circulate or reverse direction. The gradient--solenoid decomposition makes this precise: the ground state locks onto the gradient part $\grad\psi[\vek b]$, which is fixed by the flux of $\vek b$ through the boundary, whereas closed field loops inside the system are invisible to it. Since $\grad\times\vek B$ is proportional to the electric current density for static fields, fields in current-free regions are conservative and thus ideal for a CSL.

\item \emph{Baryon number.} In the chiral limit, the EoM and boundary condition imply $\int_V\vek b\cdot\grad\phi_0=\int_V(\grad\phi_0)^2$. Hence the total baryon number of the ground state is
\begin{equation*}
N_\text{B}=\frac{f_\pi^2}{\mu}\int_V(\grad\phi_0)^2=-\frac{2E[\phi_0]}{\mu}\geq0
\end{equation*}
[new Eq.~(22)]. Thus, any reduction of the condensation energy due to the inhomogeneity is accompanied by the same relative reduction of the baryon number. The local baryon number density, on the other hand, can be negative.

\item \emph{Fig.~3 (Gaussian field).} The field lines are circles centered at the origin. On a disk, the field would be tangential to the boundary everywhere, and the ground state would be exactly the trivial vacuum. On the square, the CSL is driven only by the boundary flux, which is suppressed by $e^{-L^2/(4R^2)}\approx2\times10^{-3}$. Accordingly, the baryon number density is tiny: its maximum is of order $10^{-4}$ times $\max\vek b^2$, the density that a uniform field of the same strength would induce. The pattern can be derived analytically. The boundary data are invariant under $90^\circ$ rotations and odd under reflections across the axes, which selects the leading harmonic $\phi_0\propto xz(x^2-z^2)\propto r^4\sin4\theta$. This gives $\vek b\cdot\grad\phi_0\propto-r^4e^{-r^2/R^2}\cos4\theta$: negative on the axes and positive on the diagonals, with extrema at $r=\sqrt2R\approx2.8$, exactly as in Fig.~3. The total baryon number is positive but quadratic in the small boundary flux. The lesson is that magnetic fields whose lines close inside the system are very inefficient at inducing a CSL.

\item \emph{Fig.~4 (domain wall).} The pion field would like to behave as $\pm b_0z$ on the two sides of the wall. These profiles cannot be joined, and the odd solution vanishes on the wall, and so does the baryon density. The maximum principle applied to the harmonic function $\partial_z\phi_\text{DW}$ shows that $0\leq\vek b\cdot\grad\phi_\text{DW}\leq b_0^2$, with the uniform-field value reached only on the boundaries $z=\pm L_z/2$. The less obvious lesson concerns the length scale. In the chiral limit there is no intrinsic scale, and the $L_x\to\infty$ form of the solution shows that the disturbance decays as $e^{-\pi|x|/L_z}$. The healing length $L_z/\pi$ is thus set by the size of the system along the field, not by microscopic physics. This explains why on a square domain the condensate never recovers the uniform-field profile, and why the energy excess of the wall scales as $L_z^2$ (Appendix~A). For fields well above the CSL threshold, the pion mass is only a small perturbation, and the same picture persists, as seen in Fig.~6, where the sharp-wall result stays close to the chiral-limit value $1/2$ even for $m_\pi L\gg1$.
\end{enumerate}

We have verified the statements in (b)--(d) by an independent numerical solution of the chiral-limit problem on the same domains: the identity in (b), the sign pattern, peak radius and magnitude in (c), and the bounds and healing behavior in (d).

\changes New paragraphs at the end of Sec.~3.1 (frustration mechanism, new Eq.~(22) for the baryon number). A new paragraph after the definition of the Gaussian field in Sec.~4.2, with a remark in the caption of Fig.~3 that the overall magnitude is strongly suppressed; we now refer to this configuration as the ``ground state'' rather than a ``CSL solution''. A new paragraph after the introduction of Fig.~4 in Sec.~4.3. A new paragraph in Sec.~5 that collects these observations into rules of thumb for the qualitative behavior of the ground state in more general magnetic fields.

%--------------------------------------------------------------------
\section*{Point (4): physical scales}

\begin{referee}
The authors mention in the introduction that inhomogeneities in $B$ are present in heavy-ion collisions and in neutron stars but that at least in the latter variations are on a scale much larger than relevant for modification of the original CSL state. It might be helpful to give some numbers for the scales involved to get an idea how far (if at all) the configurations seen here are away from what we would expect in physical systems. Also, would these estimates change if neutron star mergers are considered, where there are magneto-rotational instabilities and thus small-scale variations of the magnetic field? Presumably even these scales are much larger than what is needed for the configurations considered here, and also temperature effects in mergers would destroy any CSL-type phase?
\end{referee}

\reply We agree that numbers are useful, and we have added a paragraph with estimates to the introduction. In summary:
\begin{itemize}[nosep]
\item The relevant microscopic scale is $1/m_\pi\approx1.4$~fm, the width of a soliton; the CSL period is of the same order or longer. The inhomogeneities studied in our paper vary on this scale.
\item In heavy-ion collisions, $eB\sim m_\pi^2$ at top RHIC energy and about an order of magnitude more at the LHC, varying over a few fm in the transverse plane. Spatially, this is precisely the regime of our paper. However, the field decays within $\lesssim1$~fm$/c$, the baryon chemical potential is small where the field is strongest, and the temperature is far too high for the CSL order to survive.
\item In neutron stars, the chemical potential is large, but magnetar surface fields are $10^{14}$--$10^{15}$~G, and interior fields are believed not to exceed $\sim10^{18}$~G. Both are below the critical field of $\sim10^{19}$~G for $\mu\approx1$~GeV. The fields vary over macroscopic scales of up to $\sim10$~km~$\approx10^{19}$~fm.
\item In binary neutron star mergers, the Kelvin--Helmholtz and magnetorotational instabilities amplify the field to $\gtrsim10^{16}$~G and create small-scale turbulent structure. As the referee anticipated, however, the relevant scales remain macroscopic. (Simulations resolve scales of order 10~m.) On such scales, the CSL would simply follow the local field adiabatically. Moreover, merger remnants reach temperatures of several tens of MeV, which would likely destroy any CSL-type order.
\end{itemize}
The configurations considered in our paper are thus far from what is expected in neutron stars and mergers. They are, however, directly relevant for lattice QCD simulations with nonuniform (e.g.\ domain-wall-type) magnetic fields, and they serve as a theoretical laboratory for the mechanisms discussed under point~(3).

\changes A new paragraph in the introduction with the above estimates and references to the literature on magnetic fields in heavy-ion collisions [Kharzeev et al.\ 2008; Skokov et al.\ 2009; Deng and Huang 2012], magnetars [Kaspi and Beloborodov 2017], field amplification in mergers [Price and Rosswog 2006; Kiuchi et al.\ 2015], thermodynamic conditions in mergers [Perego et al.\ 2019], and thermal effects on the CSL (references already present in the paper). A sentence on the adiabatic regime of slowly varying fields has also been added to the Summary.

%--------------------------------------------------------------------
\section*{Point (5): Maxwell's equations}

\begin{referee}
Why does the EoM for $\phi$ not have to be coupled with Maxwell's equations? I understand that the coupling of the pions to $B$ is anomalous, so presumably Maxwell's equations are trivially fulfilled? Even if that's the case maybe a short remark would be helpful.
\end{referee}

\reply The referee is right. We treat $\vek B$ as an external field produced by sources outside the pion system. Varying the anomalous term $\frac{\mu}{4\pi^2}\vek B\cdot\grad\phi$ with respect to the vector potential gives the induced current
\begin{equation*}
\vek j=\frac{1}{4\pi^2}\grad\times(\mu\grad\phi)=\frac{1}{4\pi^2}\grad\mu\times\grad\phi,
\end{equation*}
which vanishes identically for uniform $\mu$, regardless of the profile of $\vek B$ and $\phi$. Equivalently, the pion field carries the magnetization $\vek M=\mu\grad\phi/(4\pi^2)$, which is curl-free. Ampère's law is therefore trivially satisfied in the bulk. The only feedback of the pion field on $\vek B$ comes from the effective magnetic charges associated with $\grad\cdot\vek M$ and with the normal component of $\vek M$ on the boundary. Upon restoring the elementary charge, the resulting correction to $\vek B$ is suppressed by $\alpha\mu^2/(4\pi^3f_\pi^2)$, which is below one percent for $\mu\approx1$~GeV. We neglect it. For a nonuniform chemical potential (Sec.~3.3), the current no longer vanishes, but its back-reaction is suppressed by $\alpha$ in the same way.

\changes A new paragraph with Eq.~(4) after the definition of $\vek b$ in Sec.~2, and a sentence in Sec.~3.3 referring to it.

%--------------------------------------------------------------------
\section*{Other changes}

\begin{itemize}[nosep]
\item Corrected ``the convolution of the coordinate dependence'' to ``the product'' in Sec.~3.3, since $\vek b$ is the product of $\mu$ and $\vek B$.
\item Corrected ``We leave both of the above equations to the future work'' to ``\dots\ questions to future work'' in Sec.~5.
\end{itemize}

\bigskip
\noindent We hope that with these changes, the manuscript is suitable for publication.

\medskip
\noindent Sincerely,\\
Tom\'a\v{s} Brauner and Ramkumar Radhakrishnan

\end{document}
